\documentclass[../../main.tex]{subfiles}
\begin{document}

{\bf [解]}
5点$A,B,C,D,E$を，$A$を原点，$AB$を$x$軸正方向にとって下図のように配置する．
各点の座標を$A(x_1,y_1), B(x_2,y_2), C(x_3,y_3), D(x_4,y_4), E(x_5,y_5)$とする．
$AB=BC=CD=DE=1$，各頂点での外角を$\theta$とする．
この時各点の座標は
\begin{align}
\begin{split}
& A(0,0),\ B(1,0),\ C(1+\cos\theta,\sin\theta), \\
& D(1+\cos\theta+\cos2\theta,\ \sin\theta+\sin2\theta), \\
& E(1+\cos\theta+\cos2\theta+\cos3\theta,\ \sin\theta+\sin2\theta+\sin3\theta)
\end{split} \label{eq:1}
\end{align}
とおける．

\begin{figure}[htb]
\begin{center}
\begin{tikzpicture}[scale=2.2, >=stealth]

  % --- パラメータ設定 (例: theta = 45度 = pi/4 の最大時) ---
  \def\thetaVal{45}

  % --- 1. 各頂点の座標計算 ---
  % A(0,0)
  \coordinate (A) at (0,0);
  % B(1,0)
  \coordinate (B) at (1,0);
  % C = B + (cos(theta), sin(theta))
  \coordinate (C) at ($(B) + ({\thetaVal}:1)$);
  % D = C + (cos(2*theta), sin(2*theta))
  \coordinate (D) at ($(C) + ({2*\thetaVal}:1)$);
  % E = D + (cos(3*theta), sin(3*theta))
  \coordinate (E) at ($(D) + ({3*\thetaVal}:1)$);

  % --- 2. 面積分割の背景ハッチング (ABC, ACD, ADE) ---
  % \triangle ABC (青系ハッチング)
  \begin{scope}
    \clip (A) -- (B) -- (C) -- cycle;
    \fill[pattern=north east lines, pattern color=blue!60] (-0.5,-0.5) rectangle (2.5,2.5);
  \end{scope}

  % \triangle ACD (赤系ハッチング)
  \begin{scope}
    \clip (A) -- (C) -- (D) -- cycle;
    \fill[pattern=north west lines, pattern color=red!60] (-0.5,-0.5) rectangle (2.5,2.5);
  \end{scope}

  % \triangle ADE (緑系ハッチング)
  \begin{scope}
    \clip (A) -- (D) -- (E) -- cycle;
    \fill[pattern=horizontal lines, pattern color=green!60!black] (-0.5,-0.5) rectangle (2.5,2.5);
  \end{scope}

  % --- 3. 五角形の辺（外枠） ---
  \draw[ultra thick, black] (A) -- (B) -- (C) -- (D) -- (E);

  % 対角線 AC, AD (分割線: 破線)
  \draw[thick, gray, dashed] (A) -- (C);
  \draw[thick, gray, dashed] (A) -- (D);
  % 閉じた五角形に見せるための EA (破線)
  \draw[thick, gray, dotted] (E) -- (A);

  % --- 4. 外角 theta を表示するための延長線 ---
  % (i) AB の延長線上での C への外角 theta
  \draw[dashed, gray] (B) -- ($(B) + (0.6, 0)$);
  \draw[->, thin, red!80!black] ($(B)+(0.3,0)$) arc (0:\thetaVal:0.3);
  \node[red!80!black, font=\small, right] at ($(B)+(0.3, 0.12)$) {$\theta$};

  % (ii) BC の延長線上での D への外角 theta
  \draw[dashed, gray] (C) -- ($(C) + ({\thetaVal}:0.6)$);
  \draw[->, thin, red!80!black] ($(C)+({\thetaVal}:0.3)$) arc (\thetaVal:{2*\thetaVal}:0.3);
  \node[red!80!black, font=\small, above right] at ($(C)+({\thetaVal+15}:0.32)$) {$\theta$};

  % (iii) CD の延長線上での E への外角 theta
  \draw[dashed, gray] (D) -- ($(D) + ({2*\thetaVal}:0.6)$);
  \draw[->, thin, red!80!black] ($(D)+({2*\thetaVal}:0.3)$) arc ({2*\thetaVal}:{3*\thetaVal}:0.3);
  \node[red!80!black, font=\small, above] at ($(D)+({2*\thetaVal+20}:0.35)$) {$\theta$};

  % --- 5. 各辺の長さラベル 1 ---
  \node[below, font=\small] at (0.5, 0) {$1$};
  \node[above left, font=\small] at ($(B)!0.5!(C)$) {$1$};
  \node[above left, font=\small] at ($(C)!0.5!(D)$) {$1$};
  \node[above right, font=\small] at ($(D)!0.5!(E)$) {$1$};

  % --- 6. 各頂点のプロットとラベル ---
  \fill[black] (A) circle (1.2pt) node[below left] {$A(0,0)$};
  \fill[black] (B) circle (1.2pt) node[below right] {$B(1,0)$};
  \fill[black] (C) circle (1.2pt) node[right] {$C$};
  \fill[black] (D) circle (1.2pt) node[above] {$D$};
  \fill[black] (E) circle (1.2pt) node[left] {$E$};

  % 面積ラベル
  \node[blue!80!black, font=\small] at (0.7, 0.2) {$\triangle ABC$};
  \node[red!80!black, font=\small] at (0.6, 0.65) {$\triangle ACD$};
  \node[green!40!black, font=\small] at (0.2, 0.85) {$\triangle ADE$};

\end{tikzpicture}
\end{center}
\caption{五角形 $ABCDE$ の配置と外角 $\theta$}
\end{figure}


五角形$ABCDE$の面積$S$は，三角形$ABC, ACD, ADE$の面積の和であって，三角形の面積公式から
\begin{align}
S &= \triangle ABC+\triangle ACD + \triangle ADE = \sum_{i=2,3,4}\dfrac{x_iy_{i+1}-x_{i+1}y_i}{2}
\end{align}
である．
ここに\cref{eq:1}を代入すると
\begin{align}
 \triangle ABC &= \dfrac{\sin\theta}{2} \\
 \triangle ACD &= \dfrac{(1+\cos\theta)(\sin\theta+\sin2\theta)-\sin\theta(1+\cos\theta+\cos2\theta)}{2} \\
               &= \dfrac{\sin2\theta+(\cos\theta\sin2\theta-\cos2\theta\sin\theta)}{2} \\
               &= \dfrac{\sin2\theta+\sin\theta}{2} \\
 \triangle ADE &= \dfrac{(1+\cos\theta+\cos2\theta)(\sin\theta+\sin2\theta+\sin3\theta)-(\sin\theta+\sin2\theta)(1+\cos\theta+\cos2\theta+\cos3\theta)}{2} \\
               &= \dfrac{(1+\cos\theta+\cos2\theta)\sin3\theta-(\sin\theta+\sin2\theta)\cos3\theta}{2} \\
               &= \dfrac{\sin3\theta + (\cos\theta\sin3\theta-\sin\theta\cos3\theta) + (\cos2\theta\sin3\theta-\sin2\theta\cos3\theta)}{2} \\ 
               & =\dfrac{\sin3\theta + \sin2\theta + \sin\theta}{2}
\end{align}
だから
\begin{align}
 S &= \dfrac{1}{2}\left[\sin\theta + \sin2\theta+\sin\theta + \sin3\theta + \sin2\theta + \sin\theta  \right] \\
   &= \dfrac{1}{2}\left[3\sin\theta + 2\sin2\theta+\sin3\theta \right] \label{eq:2}
\end{align}
と計算できる．この挙動を調べるため一階微分を計算すると
\begin{align}
 \dfrac{dS}{d\theta}
  &=\dfrac{1}{2}\left[3\cos\theta+4\cos2\theta+3\cos3\theta\right] \\
  &=\dfrac{3}{2}(\cos3\theta+\cos\theta)+2\cos2\theta \\
  &=3\cos2\theta\cos\theta+2\cos2\theta \\
  &=\cos2\theta(3\cos\theta+2)
\end{align}
となる．
$0<\theta<\pi/2$から$\cos\theta\in(0,1)$で$3\cos\theta+2>0$だから，$dS/d\theta$の符号は$\cos2\theta$と一致する．
従って$S$の増減表は以下のようになる．

\begin{table}[htb]
  \centering
\caption{$S$の増減表}
\begin{tabular}{c|ccccc}
$\theta$ & $0$ & & $\pi/4$ & & $\pi/2$ \\
\hline
$S'$ & & $+$ & $0$ & $-$ & \\
\hline
$S$ & & $\nearrow$ & & $\searrow$ &
\end{tabular}
\end{table}

従って，$S$は$\theta=\pi/4$の時，最大値
\begin{align}
S\left(\dfrac{\pi}{4}\right)=\dfrac{3}{2}\cdot\dfrac{\sqrt{2}}{2}+1+\dfrac{1}{2}\cdot\dfrac{\sqrt{2}}{2}=\sqrt{2}+1
\end{align}
をとる．


\newpage

\end{document}
